\begin{abstract}
\thispagestyle{empty}
Выпускная квалификационная работа на тему: «Разработка системы для проксирования и мониторинга имитационных сервисов в среде нагрузочного тестирования».

Цель ВКР -- проектирование и разработка автоматизированной информационной системы для централизованного управления имитационными сервисами в среде нагрузочного тестирования и мониторинга их состояния.

Метод исследования: сравнительный анализ существующих решений, объектно-ориентированное проектирование с применением нотации C4 Model, экспериментальное нагрузочное тестирование разработанного комплекса.

В первой главе проведён анализ предметной области, рассмотрены аналоги и обоснован технологический стек: Python, FastAPI, Redis. Во второй главе описано проектирование архитектуры системы, модулей проксирования, Rate Limiting и веб-интерфейса. В третьей главе рассмотрено развёртывание в ОС Astra Linux и приведены результаты функционального и нагрузочного тестирования.

Практическим результатом является программный комплекс, реализующий динамическое HTTP-проксирование (HyperText Transfer Protocol, протокол передачи гипертекста), алгоритм Fixed Window Counter с точностью ограничения запросов до 0,13~\% от теоретического значения, сквозную задержку не более 3,01 мс, а также экспорт метрик в формате Prometheus. Комплекс обеспечивает агентно-независимое управление удалёнными хостами по SSH (Secure Shell) без установки дополнительного ПО.

Библиографический список состоит из 51 источника, из них 18 печатных изданий (книги и диссертация), 5 нормативных документов (ГОСТ~-- Государственный стандарт), 28 интернет-источников.

Разработанная система позволяет отказаться от зарубежных решений в задачах управления инфраструктурой заглушек и интегрируется со стандартными средствами мониторинга.
\end{abstract}